\subsection*{4.4 Case II -- There Are No Heavy Buckets}
In this subsection, we prove that a single pseudorandom restriction preserves the expectation in the case where there is no such a heavy $B_q$. Namely, for every $q\in[Q]$, either $m_q\le3^{c^q}$ or $m_q\le\log^C n$ (or both).
\begin{authorclaim}{Lemma15}
Let $f$ be a read-once PARITY $\circ$ AND formula over $n$ variables in which the width of every term is at most $C\log n$, and in which there are no heavy buckets as described above. Then, with probability at least $1-\frac12\cdot2^{-C\log n/\log\log n}$ over $z\sim Z$ it holds that
\[
|\mathbb E[\widetilde f_z(Y)]-\mathbb E[\mathfrak f]|
\le\frac12\cdot2^{-C\log n/\log\log n}.
\]
\end{authorclaim}
Toward proving Lemma15, we partition the $Q$ buckets into two sets and treat terms that fall into each set of buckets separately. Namely, define the two sets as follows. $\mathcal A=\{q\in[Q]:m_q\le\log^{2C}n\}$. We refer to these buckets as the sparse buckets. $\mathcal B=[Q]\setminus\mathcal A$. We refer to these buckets as the well-behaved buckets.
For each set $T\in\{\mathcal A,\mathcal B\}$ we denote $f_T=\bigoplus_{i\in T}F_i$, and so $f=f_{\mathcal A}\oplus f_{\mathcal B}$. The next two subsections will be devoted to proving that the expectation of each $f_T$ is preserved after a single pseudorandom restriction. In Section4.4.3 we will combine the two results using the XOR lemma for small-bias distributions (Lemma6) to prove Lemma15.
\subsection*{4.4.1 Handling Sparse Buckets}
For the sparse buckets, we will follow the Forbes-Kelley approach [14] to prove the following.
\begin{authorclaim}{Lemma16}
With probability at least $1-\frac14\cdot2^{-C\log n/\log\log n}$ over $z\sim Z$, it holds that
\[
|\mathbb E[(\widetilde{f_{\mathcal A}})_z(Y)]-\mathbb E[\mathfrak f_{\mathcal A}]|
\le\frac14\cdot2^{-C\log n/\log\log n}.
\]
\end{authorclaim}
As outlined in Section1.5.1, Lemma16 follows readily from the work by Forbes and Kelley [14]. We require our restriction to work with high probability over $z\sim Z$, not merely in expectation, so we must redo some of Forbes and Kelley's analysis. (No substantial modification is needed.) The details follow.
\begin{proof}[Proof of Lemma16]
First, recall that each term in $f_{\mathcal A}$ is of width at most $C\log n$. There are at most $\log^{2C}n$ terms in each bucket, and at most $Q=O(\log\log n)$ such buckets, so overall $f_{\mathcal A}$ reads at most $n'=\log^{2C+2}n$ variables.
Note that $f_{\mathcal A}$ can be computed by a width-4 ROBP of length $n'$. We follow [14] and let $G:\{0,1\}^{n'}\to\mathbb R^{4\times4}$ encode the transition of the branching program. Namely, perhaps after renumbering the variables, we have $G(x)=G_1(x_1)\cdots G_{n'}(x_{n'})$ where $G_i(x_i)=A_{i,x_i}$ for $A_{i,b}$ being the transition matrix that corresponds to taking the bit $b$ while at layer $i$. Set $k_0=8\log n/\log\log n$, and note that $k_0\le k$. By [14, Lemma4.1], $G$ can be written as
\[
G=\mathbb E[G]+L+\sum_{i=1}^{n'}H_i\cdot G^{>i},
\]
where $L$ has degree\footnote{We say a function $H:\{0,1\}^n\to\mathbb R^{w\times w}$ having Fourier expansion $\sum_{\alpha\in\mathbb F_2^n}\widehat H_\alpha\chi_\alpha$ has degree $d$ if $\widehat H_\alpha$ is the zero matrix for every $\alpha$ with Hamming weight larger than $d$.} less than $k_0$, $H_i$ is of degree exactly $k_0$, $G^{>i}$ is a width-4 ROBP, and $H_i$ and $G^{>i}$ are on disjoint set of variables. More specifically,
\[
L=\sum_{\substack{\alpha\in\mathbb F_2^{n'}\\0<|\alpha|<k_0}}\widehat G_\alpha\chi_\alpha
\]
is the truncated Fourier expansion of $G$, $G^{>i}(x_{i+1},\ldots,x_n)=G_{i+1}(x_{i+1})\cdots G_{n'}(x_{n'})$, and
\[
H_i=\sum_{\substack{\alpha\in\mathbb F_2^{n'}\\|\alpha|=k_0,\alpha_i=1}}\widehat{G^{\le i}}_\alpha\chi_\alpha,
\]
where $G^{\le i}(x_1,\ldots,x_i)=G_1(x_1)\cdots G_i(x_i)$. Let $\|\cdot\|$ be the Frobenius norm. By sub-additivity, we have
\begin{align}
\mathbb E_Z[\|\mathbb E_{Y,U}[G(Y+Z\wedge U)]-\mathbb E[G]\|]
&\le\mathbb E_Z[\|\mathbb E_{Y,U}[L(Y+Z\wedge U)]\|]\notag\\
&\quad+\sum_{i=1}^{n'}\mathbb E_Z[\|\mathbb E_{Y,U}[(H_i\cdot G^{>i})(Y+Z\wedge U)]\|].\tag{7}
\end{align}
Just as in [14], the low-degree term $L$ is dealt with a $\delta$-biased distribution. From the work of Chattopadhyay, Hatami, Reingold, and Tal [9] we know that
\[
L_1(L)=\sum_{k'_0=1}^{k_0}(c_{\mathrm{CHRT}}\log n')^{4k'_0}\le2(c_{\mathrm{CHRT}}\log n')^{4k_0}
\]
for some universal constant $c_{\mathrm{CHRT}}\ge1$. Thus, by Claim5, we get that the first term of Equation(7) is bounded by
\[
2\delta\cdot2(c_{\mathrm{CHRT}}\log n')^{4k_0}
\le2^{-C\log n}\cdot2^{8k_0\log\log\log n}\le n^{-C/2},
\]
taking into account the fact that $\mathbb E[L(U)]=0$.
For each $i$ of the second term of Equation(7), we use sub-multiplicativity and the fact that $H_i$ and $G^{>i}$ are on disjoint set of variables to get
\[
\mathbb E_Z[\|\mathbb E_{Y,U}[(H_i\cdot G^{>i})(Y+Z\wedge U)]\|]
\le\mathbb E_{Y,Z}[\|\mathbb E_U[H_i(Y+Z\wedge U)]\|\cdot\|\mathbb E_U[G^{>i}(Y+Z\wedge U)]\|].
\]
As $G^{>i}$ is a width-4 ROBP, $\|\mathbb E_U[G^{>i}(y+z\wedge U)]\|\le2$ for all $y\sim Y$ and $z\sim Z$. Continuing the above bound, by Cauchy-Schwarz we get
\[
\mathbb E_Z[\|\mathbb E_{Y,U}[(H_i\cdot G^{>i})[Y+Z\wedge U]]\|]
\le2\sqrt{\mathbb E_{Y,Z}[\|\mathbb E_U[H_i(Y+Z\wedge U)]\|^2]}.
\]
Following [14, Lemma7.1]\footnote{Forbes and Kelley [14] take the bits of $Z$ to have marginals $p=1/2$, but one can extend the lemma easily for the case of a general $p$.}, using the bound by Chattopadhyay et al. [9] and Parseval's identity [14, Proposition3.1], we get
\begin{align*}
\mathbb E_{Y,Z}[\|\mathbb E_U[H_i(Y+Z\wedge U)]\|^2]
&\le(2^{-Ck_0}+\gamma)\cdot\left(\delta\left(\sum_{\alpha\in\mathbb F_2^{n'}}\|(\widehat H_i)_\alpha\|\right)^2+\sum_{\alpha\in\mathbb F_2^{n'}}\|(\widehat H_i)_\alpha\|^2\right)\\
&\le(2^{-Ck_0}+\gamma)\cdot\left(\delta\cdot L_1^2(G^{\le i})+\mathbb E[\|G^{\le i}(U)\|^2]\right)
\le8\cdot2^{-Ck_0}.
\end{align*}
Overall, we get that
\begin{align*}
\mathbb E_Z[\|\mathbb E_{Y,U}[G(Y+Z\wedge U)]-\mathbb E[G]\|]
&\le n^{-C/2}+2n'\sqrt{8\cdot2^{-Ck_0}}\\
&\le\frac1{16}\cdot2^{-Ck_0/4}
=\frac1{16}\cdot2^{-2C\log n/\log\log n},
\end{align*}
and we can choose the encoding $G$ so that $f_{\mathcal A}(x)=G(x)_{1,1}$. Markov's inequality completes the proof.
\end{proof}
